\section{Project Synthesis}

\begin{studybox}{Project synthesis: what makes the 1962 calendar the kind of thing it is}
This block states an editorial judgement of the project, with its reasons, the strongest
counterargument at full strength, and the specific finding that would change it. It is a
source-grounded synthesis, not a claim about the intention of any legislator and not a
theological or juridical determination.
\end{studybox}

\subsection{The judgement}

The decisive innovation of the 1960 code is not the replacement of the double and semidouble
grades by four classes. It is the exclusivity clause: \rn{7}'s statement that precedence among
liturgical days \latin{unice per peculiarem tabellam determinatur}, and \rn{91}'s repetition
that the table governs \latin{sublatis quibuslibet aliis titulis vel normis}. What 1960 did was
to move the entire question of precedence out of a distributed body of grades, privileges,
customs, and responses to doubts, and into one enumerated table of twenty-eight lines — and then
to revoke, at \rn{3} of \latin{Rubricarum instructum}, everything that stood against it,
including immemorial custom and privileges deserving individual mention. The four classes are a
consequence of that move, not its substance: they are names for regions of the table.

The second half of the judgement is that the resumption mechanism of \rn{18} is the clearest
surviving evidence that the 1960 reform was a codification and not a rationalisation. Every
other structural feature of the calendar was made decidable: precedence by table, translation by
\rn{96}, reposition by \rn{101}, commemorations by a count. \rn{18} alone retains a rule that is
not decidable from its own terms — it enumerates four cases and leaves a fifth, the
twenty-three-Sunday year, to whatever the closing \latin{semper} clause is taken to mean; and it
takes its wording, essentially unchanged, from the books it replaced. Where the reform could
codify, it codified sharply. Where the existing rule already worked in the ordinary case, it was
carried over as it stood, gap and all.

The reasons are three, and each rests on evidence set out above.

\begin{enumerate}[leftmargin=1.6em, itemsep=0.3em]
\item \textbf{The code says so, twice, in the strongest available terms.} Exclusivity is
asserted at \rn{7} and again at \rn{91}, and the revocation clause of \latin{Rubricarum
instructum} \rn{3} is drafted to defeat exactly the kind of local claim that a distributed
system of dignities generates. A reform whose point was the renaming of grades would not need
that clause.
\item \textbf{The table does work the classes cannot do.} A II class Sunday and a II class feast
of the Lord are of the same class and are separated by lines 14 and 15; a Lenten feria and an
Advent feria are of the same class and are separated by lines 22 and 25; a particular III class
feast outranks a universal one at lines 23 and 24. In each case the class label is silent and
the table decides. If the classes carried the reform, these distinctions would have had to be
expressed as further classes, and they are not.
\item \textbf{The one rule the reform did not make decidable is the one it did not rewrite.}
The resumption rubric printed at the end of the Twenty-third Sunday after Pentecost in the 1962
Missal is, word for word, the rubric printed in the same place in the identified pre-1955
witness, differing only in its internal page cross-reference and in the grade named for the
resumed formularies. \rn{18} of the code restates the same four
cases. Its enumeration begins at twenty-five and stops at twenty-eight, and the shortfall case
that the Missal's own perpetual table shows to be possible is nowhere addressed. That is the
signature of carried-over text, not of fresh drafting.
\end{enumerate}

\subsection{The strongest counterargument}

The counterargument, stated at full strength, runs as follows.

Exclusivity clauses of this kind are not innovations at all; they are the ordinary drafting
furniture of a code, and the pre-1960 books had their own tables of occurrence and concurrence
serving the same function. What actually changed in 1960 was the substantive content of the
calendar: an entire grade, the semidouble, disappeared into the III class; a whole stratum of
simples became commemorations; every octave but three was abolished; every vigil but seven was
abolished; the number of orations admitted at a Mass was capped at three by Missal rubrics
\rn{435} with the emphatic \latin{nullo praetextu}; the Suffrage of All Saints and the
commemoration of the Cross were abolished outright; and \rn{14} destroyed the resumption of
impeded Sundays on weekdays. Those are changes a worshipper would notice on an ordinary Tuesday.
The table, by contrast, would have produced the same result on most days under either drafting.
To say that the exclusivity clause is the decisive innovation is to mistake the form of the code
for its effect, and to elevate a drafting technique above the substantive suppressions that the
\latin{Variationes} record item by item.

And on the second half: \rn{18} is not evidence of a codifying temper. It is simply an
inherited rule that nobody noticed had a gap, because the gap is engaged in four years out of
two hundred. Reading a philosophy of reform out of a rubric that fails once every half century
is over-reading. A better explanation is banal: the drafters enumerated the cases that occur
often enough to matter and did not audit the arithmetic at the tails.

\subsection{Reply, and what would change the judgement}

The counterargument is right that the substantive suppressions are what a worshipper notices,
and this reference has set them out at their loci precisely because they are the visible content
of the reform. The judgement above is not that the suppressions are unimportant; it is about
what makes the resulting calendar the kind of object it is — an object in which a stated
question about precedence has, in principle, exactly one answer derivable from one table. That
is a property of the 1962 calendar that the pre-1960 calendar did not have to the same degree,
and it is the property on which every downstream reference, Ordo, and computation depends.

The reply to the second point is narrower. That a gap engages rarely does not make it
uninformative; what makes it informative is that the surrounding rules were rewritten and this
one was not, and that the book carrying the rule also prints, in its own front matter, the perpetual
table showing that the unlegislated case arises. The evidence of carried-over text is textual,
not statistical.

\begin{studybox}{What would change this judgement}
\textbf{On the first half.} A demonstration that the pre-1960 general rubrics already contained
an equivalent exclusivity clause — that is, an identified locus in the pre-1955 \latin{Rubricae
generales Breviarii et Missalis Romani} or in the \latin{Additiones et Variationes} under
\latin{Divino afflatu} stating that precedence is determined solely by a table and by nothing
else — would remove the contrast on which the judgement rests. Half of that test has now been
run and the judgement survives it: the identified Benziger \latin{editio IV iuxta typicam
Vaticanam} described in Section~\ref{sec:witness} carries the pre-1955 \latin{Rubricae generales
Missalis} and the \latin{Additiones et Variationes} in full, and they contain no table of
precedence and no exclusivity clause; they settle occurrence by the comparative
\latin{Officium nobilius}, which is precisely the distributed idiom of dignities that \rn{7} and
\rn{91} displace. The other half is untouched. The pre-1955 tables of occurrence and concurrence
were printed in the Breviary, not the Missal, and no pre-1955 Breviary was read for this
edition; if such a table were found to be introduced by an equivalent exclusivity clause, the
contrast would be a contrast of books rather than of systems, and the judgement would fail. That
remains the single most decisive available test.

\textbf{On the second half.} Evidence from the preparatory work of the commission for the
general liturgical restoration — a schema, a \latin{votum}, or a recorded discussion — showing
that the resumption mechanism was examined and deliberately retained, or that the shortfall case
was raised and answered, would convert \rn{18} from carried-over text into a considered
decision, and the inference from it to a codifying rather than rationalising temper would fail.
Conversely, a competent Ordo for 1943, 2011, 2038, or 2095 resolving the shortfall case would
show that the practice had an answer even though the code did not, which would weaken but not
destroy the point about the text.
\end{studybox}
